% ============================================================
%  Chapitre 1 — Introduction
% ============================================================
\chapter*{Introduction}
\addcontentsline{toc}{chapter}{Introduction}
\label{chap:intro}

% ----------------------------------------------------------------
\section{Contexte et motivation}
% ----------------------------------------------------------------

Les algorithmes d'\emph{Intelligence en essaim} (\textit{Swarm Intelligence}) ont été
introduits par Beni et Wang en 1989~\cite{beni1989swarm}. Cette famille de méthodes
s'inspire du comportement collectif d'insectes sociaux : des agents simples, dotés
d'une capacité de traitement locale limitée, interagissent \emph{indirectement} via
l'environnement — un mécanisme connu sous le nom de \textbf{stigmergie} — pour produire
collectivement un comportement adaptatif complexe, sans aucune coordination centralisée.

L'algorithme d'essaim le plus étudié est l'\textbf{Optimisation par Colonies de Fourmis}
(\textit{Ant Colony Optimisation}, ACO), formalisé par Dorigo, Maniezzo et Colorni
en 1996~\cite{dorigo1996ant}. Son inspiration directe est le fourragement des fourmis
réelles : en se déplaçant, chaque fourmi dépose des \emph{phéromones} sur son trajet ;
les congénères perçoivent ces phéromones et tendent à emprunter les chemins les plus
concentrés, ce qui renforce progressivement les routes les plus courtes tout en laissant
s'évaporer les marques inutilisées. Cette rétroaction positive/négative converge vers
une solution proche de l'optimum global~\cite{dorigo2004ant}.

Dans ce projet, on met en œuvre un modèle ACO simplifié pour résoudre le
\textbf{problème de fourragement} sur un terrain fractal 2D : trouver le chemin de
coût minimal entre le nid et une source de nourriture, en tenant compte d'un terrain
à relief variable généré par un algorithme de plasma fractal (subdivision récursive
\textit{midpoint displacement}~\cite{fournier1982midpoint}).

\section{Objectifs du rapport}
% ----------------------------------------------------------------

Ce rapport suit la progression imposée par le sujet de TP~\cite{ensta2026os202} :

\begin{enumerate}
    \item \textbf{Instrumentation} : mesure expérimentale des temps par phase sur le
          code de référence (\texttt{src\_v1\_timing/}) ;
    \item \textbf{Vectorisation} : passage AoS $\rightarrow$ SoA et analyse de l'impact
          sur les performances mono-c\oe{}ur (\texttt{src\_v2\_soa/}) ;
    \item \textbf{OpenMP} : parallélisation en mémoire partagée, mesure de l'accélération
          $S(p)$ en fonction du nombre de threads (\texttt{src\_v3\_omp/}) ;
    \item \textbf{MPI} : deux stratégies de parallélisation distribuée, analyse du
          rapport calcul/communication (\texttt{src\_v4\_mpi/}).
\end{enumerate}

Les mesures sont réalisées sur deux machines distinctes (chapitre~\ref{chap:env})
afin d'évaluer la portabilité des conclusions.

% ----------------------------------------------------------------
\section{Organisation du rapport}
% ----------------------------------------------------------------

Le chapitre~\ref{chap:modele} décrit le modèle ACO et son implémentation.
Le chapitre~\ref{chap:env} présente les environnements expérimentaux (deux machines)
et le protocole de mesure. Le chapitre~\ref{chap:baseline} analyse le code de
référence et l'effet de la vectorisation. Les chapitres~\ref{chap:openmp}
et~\ref{chap:mpi} présentent respectivement les résultats OpenMP et MPI.
Le chapitre~\ref{chap:discussion} discute l'ensemble des résultats et conclut.
